\subsection{拟紧空间与紧空间}

定义 1. 拓扑空间 X 称为拟紧的，如果它满足以下公理：

\begin{enumerate}
\def\labelenumi{(\Alph{enumi})}
\setcounter{enumi}{2}
\tightlist
\item
  X 上的每个滤子都至少有一个聚点。
\end{enumerate}

拓扑空间称为紧的，如果它既是拟紧的又是豪斯多夫的。

由这个公理立即得出：若 \(f\) 是集合 \(Z\) 到拟紧空间 \(X\) 中的一个映射，而 \(\mathfrak{F}\) 是 \(Z\) 上的任一滤子，则 \(f\) 关于 \(\mathfrak{F}\) 至少有一个聚点。特别地，拟紧空间的每个点列都至少有一个聚点；但这个条件并不等价于 (C)（习题 11）。

我们给出三个公理，其中每一个都等价于公理 (C)：

(C') \(X\) 上的每个超滤子都收敛。

\((\mathrm{C}^{\prime}) \Longrightarrow (\mathrm{C})\)：若 \(\mathfrak{F}\) 是 \(\mathbf{X}\) 上的一个滤子，则存在比 \(\mathfrak{F}\) 更细的超滤子（§ 6，第 4 目，定理 1）。既然这个超滤子收敛于一点 \(x\)，\(x\) 就是 \(\mathfrak{F}\) 的聚点。

\begin{enumerate}
\def\labelenumi{(\Alph{enumi})}
\setcounter{enumi}{2}
\tightlist
\item
  \(\Longrightarrow\) (C')：因为若一个超滤子有一个聚点，则它收敛于该点（§ 7，第 2 目，命题 4 的推论）。
\end{enumerate}

若 \(f\) 是集合 \(Z\) 到拟紧空间 \(X\) 中的一个映射，而 \(\mathfrak{U}\) 是 \(Z\) 上的一个超滤子，则 \(f\) 关于 \(\mathfrak{U}\) 至少有一个极限点（§ 6，第 6 目，命题 10）。

(C'') X 的闭子集的每个交为空的族都包含一个交为空的有限子族。

\begin{enumerate}
\def\labelenumi{(\Alph{enumi})}
\setcounter{enumi}{2}
\tightlist
\item
  \(\Longrightarrow\) (C'\,')：设 \(\mathfrak{G}\) 是 X 的闭子集的一个交为空的族。若 \(\mathfrak{G}\) 的每个有限子族都有非空的交，则 \(\mathfrak{G}\) 生成一个滤子（§ 6，第 2 目，命题 1），由假设它有一个聚点。这个点属于 \(\mathfrak{G}\) 的所有集合（因为它们是闭的）；于是得到矛盾。
\end{enumerate}

\((\mathbf{C}^{\prime\prime}) \Longrightarrow (\mathbf{C}) : \text{因为若 } (\mathbf{C}) \text{为假，则存在一个滤子 } \mathfrak{F} \text{在 } \mathbf{X} \text{上没有聚点；因此滤子 } \mathfrak{F} \text{中的各集合的闭包构成 } \mathbf{X} \text{的一个闭子集族，违背公理 } (\mathbf{C}^{\prime\prime}).\)

(C \(^{iii}\) )（Borel-Lebesgue 公理）X 的每个开覆盖都包含 X 的一个有限开覆盖。

\((\mathbf{C}^{\prime\prime\prime}) \iff (\mathbf{C}^{\prime\prime})\)，取余集即得。

若 X 是拟紧的，则 X 的每个局部有限覆盖 R 都是有限的。因为 X 的每一点都有一个只与 R 中有限多个集合相交的开邻域，而由 \((\mathbf{C}^{\prime\prime})\)，这些邻域中的有限多个覆盖 X。

例. 1) 每个有限空间都是拟紧的，更一般地，每个只含有限多个开集的空间都是拟紧的。有限空间是紧的当且仅当它是离散的，因为有限豪斯多夫空间是离散的（§ 8，第 1 目，命题 3 的推论）。反之，每个紧的离散空间都是有限的，因为在这种空间中由单点组成的诸集合都是开集；于是由 (C'\,')，该空间是有限的。

\begin{enumerate}
\def\labelenumi{\arabic{enumi})}
\setcounter{enumi}{1}
\tightlist
\item
  设 X 是一个集合，给 X 赋以这样的拓扑：其闭集为 X 及 X 的所有有限子集{[}这个子集族显然满足 §1 第 4 目的公理 \((\mathrm{O}_{\mathrm{I}}^{\prime})\) 与 \((\mathrm{O}_{\mathrm{II}}^{\prime})\){]}。这样定义的拓扑空间是拟紧的。因为若 \((\mathbf{F}_{i})_{i\in I}\) 是 X 的闭子集的一个交为空的族，则 \(F_{\alpha}\) 对某个 \(\alpha\in I\) 是有限的。设 \(a_{k}\)（\(1\leqslant k\leqslant n\)）是 \(F_{\alpha}\) 的诸元素；则由假设，对每个指标 k 存在指标 \(i_{k}\in I\) 使 \(a_{k}\notin F_{i_{k}}\)；于是诸 \(F_{i_{k}}\)（\(1\leqslant k\leqslant n\)）与 \(F_{\alpha}\) 的交为空，从而公理 \((\mathrm{C}^{\prime\prime})\) 得以满足。若 X 是无限集，它就不是豪斯多夫的。
\end{enumerate}

注. 拟紧（非豪斯多夫）空间主要用于拓扑学对代数几何的应用，很少出现在其他数学理论中；与之相反，紧空间在其他数学理论中扮演着重要角色。

定理 1. 设 \(\mathfrak{F}\) 是拟紧空间 \(\mathbf{X}\) 上的一个滤子，并设 \(\mathbf{A}\) 是 \(\mathfrak{F}\) 的聚点组成的集合。则 \(\mathbf{A}\) 的每个邻域都属于 \(\mathfrak{F}\)。

设 V 是 A 的一个邻域，并假设 F 的每个集合都与 CV 相交。于是 F 的诸集合与 CV 的交构成 X 上一个滤子 G 的基；X 是拟紧的，故 G 至少有一个聚点 y，而 y 不属于 A，因为 A 的邻域 V 不与 G 的某些集合相交。但既然 G 比 F 更细，y 也是 F 的聚点，这与假设矛盾。

推论. 紧空间上的一个滤子收敛，当且仅当它只有一个聚点。

必要性由 § 8，第 1 目，命题 1 得出；充分性由上面的定理 1 得出。

